\section{Main Results}\label{sec:main-results}
We state the eight no-go theorems. Each theorem is followed by a brief remark clarifying scope.

\begin{theorem}[NG\_ARROW\_DPI]\label{thm:NG_ARROW_DPI}
\ThmStmtNGARROWDPI
\end{theorem}

\begin{remark}
\textbf{Interpretation.} This theorem rules out any increase of forward--reverse path-space KL under an honest deterministic observation. In the finite audited setting, coarse observation cannot manufacture a larger arrow than the micro path law already carries.

\textbf{Scope / escape routes.} The obstruction is evaded only by leaving the theorem object: stochastic observation, fitted macro chains, or other proxy reconstructions are outside this claim. Real microscopic nonreversibility can still produce positive arrow; the theorem states that deterministic coarse-graining does not amplify it.

\end{remark}


\begin{theorem}[NG\_PROTOCOL\_TRAP]\label{thm:NG_PROTOCOL_TRAP}
\ThmStmtNGPROTOCOLTRAP
\end{theorem}

\begin{remark}
\textbf{Interpretation.} This theorem rules out honest observed arrow for a stationary reversible lift, even when hidden protocol coordinates are present. Hidden protocol structure by itself does not create entropy production in the audited stationary regime.

\textbf{Scope / escape routes.} The scope changes once stationarity or reversibility is dropped. Nonstationary initialization, external scheduling, or genuinely driven lifted dynamics fall outside the theorem hypotheses and can reopen positive observed arrow.

\end{remark}


\begin{theorem}[NG\_FORCE\_FOREST]\label{thm:NG_FORCE_FOREST}
\ThmStmtNGFORCEFOREST
\end{theorem}

\begin{remark}
\textbf{Interpretation.} This theorem rules out force-like cycle obstruction on forest support. On a finite forest, every antisymmetric edge field is potential-generated, so there is no cycle-supported drive to sustain a nontrivial affinity class.

\textbf{Scope / escape routes.} The obstruction is evaded by allowing non-forest support or by changing the exact support notion. Positive cycle rank reopens the possibility of nontrivial circulation, while thresholded proxy graphs are outside the theorem object.

\end{remark}


\begin{theorem}[NG\_FORCE\_NULL]\label{thm:NG_FORCE_NULL}
\ThmStmtNGFORCENULL
\end{theorem}

\begin{remark}
\textbf{Interpretation.} This theorem rules out cycle affinity when the log-ratio 1-form is exact. Exact ratio-form structure collapses every closed-walk sum, so the force-like class is null in the precise bidirected-support sense of the theorem.

\textbf{Scope / escape routes.} The obstruction is evaded by breaking exactness or by replacing the exact support with a thresholded or regularized proxy graph. The theorem is about exact bidirected support and exact ratio-form data, not about diagnostic pictures.

\end{remark}


\begin{theorem}[NG\_MACRO\_CLOSURE\_DEFICIT]\label{thm:NG_MACRO_CLOSURE_DEFICIT}
\ThmStmtNGMACROCLOSUREDEFICIT
\end{theorem}

\begin{remark}
\textbf{Interpretation.} This theorem identifies closure deficit as the exact obstruction to a closed macro law at the chosen lag. Positive conditional-information gap means that no single macro kernel captures the packaged future laws without loss.

\textbf{Scope / escape routes.} The obstruction is evaded only by changing the package, lag, or underlying dynamics, or by accepting proxy fitted kernels as diagnostics rather than theorem objects. A fitted macro chain may still be useful operationally, but it is not theorem evidence for closure.

\end{remark}


\begin{theorem}[NG\_OBJECT\_CONTRACTIVE]\label{thm:NG_OBJECT_CONTRACTIVE}
\ThmStmtNGOBJECTCONTRACTIVE
\end{theorem}

\begin{remark}
\textbf{Interpretation.} This theorem bounds epsilon-stable packaged states in a strict-contraction regime and yields uniqueness in the exact fixed-point case. Within the admissible TV/Dobrushin hypotheses, multiple well-separated robust objects cannot coexist at that scale.

\textbf{Scope / escape routes.} The obstruction is evaded by moving to noncontractive regimes or by using clustering heuristics that are outside theorem evidence. The claim is deliberately bounded to the contraction setup and does not attempt a general objecthood theory.

\end{remark}


\begin{theorem}[NG\_LADDER\_IDEM]\label{thm:NG_LADDER_IDEM}
\ThmStmtNGLADDERIDEM
\end{theorem}

\begin{remark}
\textbf{Interpretation.} This theorem rules out any iterate ladder from an idempotent packaging map. Once the operator is applied once, later iterates add no new structure because the map has already stabilized.

\textbf{Scope / escape routes.} The obstruction is evaded only by changing the operator class or by changing the package itself. Non-idempotent updates, interface growth, or theory changes lie outside the theorem hypotheses.

\end{remark}


\begin{theorem}[NG\_LADDER\_BOUNDED\_INTERFACE]\label{thm:NG_LADDER_BOUNDED_INTERFACE}
\ThmStmtNGLADDERBOUNDEDINTERFACE
\end{theorem}

\begin{remark}
\textbf{Interpretation.} This theorem rules out an infinite strictly increasing ladder of predicates once the interface is fixed and finite. The finite image of the lens gives a finite algebra of definable predicates, so internal ladder growth must stabilize.

\textbf{Scope / escape routes.} The obstruction is evaded only by changing the theory: lens growth, domain growth, or package change can restart definable expansion. With a fixed finite interface, no unbounded internal ladder remains.

\end{remark}


\begin{remark}[Supplementary apparatus]
Supplementary material collects the reproducibility map, including computational audit artifacts and a formalization-status table. The statements and proofs in the main text do not depend on those materials.
\end{remark}
